\section{Clipped-KL descent in split coordinates}
\label{app:nonconvex-details}

We prove the one-step estimate in
\Cref{thm:nonconvex-clipped-stationarity}.  The point that requires care is
the orientation of the Gaussian divergences.  The forward divergence controls
the smoothness remainder and appears in the proximal model, whereas the
covariance part of the reverse divergence is what remains in the descent
certificate.

Write
\[
  \calE(m,C)=\Phi(C)+\calE_2(m,C)+\mathrm{const},
  \qquad
  \Phi(C):=-\frac12\log\det C,
  \qquad
  \calE_2(m,C):=\E_{\N(m,C)}V(X).
\]
For two feasible states \(a=(m,C)\) and \(a'=(m',C')\), put
\begin{align*}
  D_m(a',a)
  &:=\frac12\norm{m'-m}_{C^{-1}}^2,\\
  D_C^+(C',C)
  &:=\KL\!\left(\N(0,C')\,\middle\|\,\N(0,C)\right),\\
  D_C^-(C,C')
  &:=\KL\!\left(\N(0,C)\,\middle\|\,\N(0,C')\right).
\end{align*}

\begin{lemma}[Relative smoothness on the covariance box]
\label{lem:app:nonconvex-relative-smoothness}
Suppose
\(-\beta\Id\preceq\nabla^2V\preceq\beta\Id\) and both \(C,C'\) lie in
\([\lambda_-\Id,\lambda_+\Id]\).  If
\(b=\E_{\N(m,C)}\nabla V\) and
\(A=\E_{\N(m,C)}\nabla^2V\), then
\begin{equation}
  \calE_2(m',C')
  \leq
  \calE_2(m,C)
  +b^\top(m'-m)+\frac12\Tr\!\left(A(C'-C)\right)
  +2\beta\lambda_+
   \left\{D_m(a',a)+D_C^+(C',C)\right\}.
  \label{eq:app:nonconvex-relative-smoothness}
\end{equation}
\end{lemma}

\begin{proof}
Set \(u=m'-m\),
\(M=C^{-1/2}C'C^{-1/2}\),
\(B_0=C^{1/2}\),
\(B'=C^{1/2}M^{1/2}\), and \(H=B'-B_0\).  Then
\(B'B'^{\mathsf T}=C'\).  Couple
\[
  X=m+B_0Z,
  \qquad
  Y=m'+B'Z,
  \qquad Z\sim\N(0,\Id).
\]
The upper Hessian bound in absolute value gives
\[
  \E V(Y)
  \leq
  \E V(X)
  +\E\!\left[\nabla V(X)^\top(u+HZ)\right]
  +\frac\beta2\left(\norm u^2+\norm H_{\Frob}^2\right).
\]
Gaussian integration by parts and
\[
  C'-C=B_0H^{\mathsf T}+HB_0^{\mathsf T}+HH^{\mathsf T}
\]
show that
\begin{align*}
  \E\!\left[\nabla V(X)^\top HZ\right]
  &=\Tr(H^{\mathsf T}AB_0)\\
  &=\frac12\Tr\!\left(A(C'-C)\right)
    -\frac12\Tr(AHH^{\mathsf T}).
\end{align*}
Since \(\norm A_{\op}\leq\beta\), the last term and the Taylor remainder
together are bounded above by
\[
  \frac\beta2\norm u^2+\beta\norm H_{\Frob}^2.
\]
The covariance box yields
\[
  \frac\beta2\norm u^2
  \leq\beta\lambda_+D_m(a',a).
\]
Moreover,
\(\norm H_{\Frob}^2
\leq\lambda_+\norm{M^{1/2}-\Id}_{\Frob}^2\), while the scalar inequality
\[
  x-1-\log x\geq(\sqrt x-1)^2,
  \qquad x>0,
\]
gives
\[
  \norm{M^{1/2}-\Id}_{\Frob}^2
  \leq
  \sum_i\{\lambda_i(M)-1-\log\lambda_i(M)\}
  =2D_C^+(C',C).
\]
Substitution proves \eqref{eq:app:nonconvex-relative-smoothness}.
\end{proof}

We next identify the clipped update as a constrained proximal step.  At
\(a_n=(m_n,C_n)\), define
\begin{equation}
\label{eq:app:nonconvex-proximal-model}
  \mathcal Q_n(m,C)
  :=
  \Phi(C)+b_n^\top(m-m_n)
  +\frac12\Tr\!\left(A_n(C-C_n)\right)
  +\frac1h\left\{D_m((m,C),a_n)+D_C^+(C,C_n)\right\}.
\end{equation}
With the trace pairing on symmetric matrices,
\begin{equation}
  D_C^+(C,C_n)
  =\Phi(C)-\Phi(C_n)
   -\Tr\!\left(\nabla\Phi(C_n)(C-C_n)\right),
  \qquad
  \nabla\Phi(C)=-\frac12C^{-1}.
  \label{eq:app:nonconvex-forward-is-bregman}
\end{equation}
Thus the covariance penalty in \(\mathcal Q_n\) is the
log-determinant Bregman divergence in its first argument.

The mean and covariance variables separate.  The mean minimizer is
\(m_n-hC_nb_n\).  The unconstrained covariance minimizer is
\[
  \widetilde C_{n+1}
  =(1+h)(C_n^{-1}+hA_n)^{-1}.
\]
Under
\(h\beta\lambda_+\leq1/2\), its inverse is positive definite because
\[
  C_n^{-1}+hA_n
  \succeq(\lambda_+^{-1}-h\beta)\Id
  \succeq(2\lambda_+)^{-1}\Id.
\]
More precisely, if \(\alpha=1+h^{-1}\), the covariance-dependent part of
\(\mathcal Q_n\) is, up to a constant,
\begin{equation}
  \alpha\,D_\Phi(C,\widetilde C_{n+1}),
  \qquad
  D_\Phi(C,\widetilde C)
  :=\Phi(C)-\Phi(\widetilde C)
    -\Tr\!\left(\nabla\Phi(\widetilde C)(C-\widetilde C)\right).
  \label{eq:app:nonconvex-proximal-bregman-form}
\end{equation}
Indeed,
\(\alpha\widetilde C_{n+1}^{-1}=h^{-1}C_n^{-1}+A_n\), which matches the
linear coefficient in \eqref{eq:app:nonconvex-proximal-model}.

For completeness, the Bregman projection in
\eqref{eq:app:nonconvex-proximal-bregman-form} is exactly spectral clipping.
Write
\(\widetilde C_{n+1}=Q\diag(\widetilde c_i)Q^\top\) and
\(C^+=Q\diag(c_i^+)Q^\top\), where
\(c_i^+=\min\{\lambda_+,\max\{\lambda_-,\widetilde c_i\}\}\).  Then
\[
  G:=\nabla_C D_\Phi(C^+,\widetilde C_{n+1})
  =\frac12\bigl(\widetilde C_{n+1}^{-1}-(C^+)^{-1}\bigr)
  =Q\diag(g_i)Q^\top.
\]
Here \(g_i>0\) only when \(c_i^+=\lambda_-\), \(g_i<0\) only when
\(c_i^+=\lambda_+\), and \(g_i=0\) otherwise.  Hence every feasible
\(C\) satisfies
\[
  \Tr\!\left(G(C-C^+)\right)
  =\sum_i g_i\{q_i^\top Cq_i-c_i^+\}\geq0.
\]
This is the variational inequality for minimizing
\(D_\Phi(\,\cdot\,,\widetilde C_{n+1})\) over the covariance box; strict
convexity in the first argument gives uniqueness.  Consequently,
\(C^+=\operatorname{Clip}_{[\lambda_-,\lambda_+]}
(\widetilde C_{n+1})=C_{n+1}\), and
\((m_{n+1},C_{n+1})\) is the unique minimizer of
\eqref{eq:app:nonconvex-proximal-model} over the covariance box.

For the mean minimizer, with \(u_n=m_{n+1}-m_n\),
\begin{equation}
  b_n^\top u_n=-\frac1h\norm{u_n}_{C_n^{-1}}^2
  =-\frac2hD_m(a_{n+1},a_n).
  \label{eq:app:nonconvex-mean-optimality}
\end{equation}
For the covariance minimizer, constrained first-order optimality gives, for
every feasible \(C\),
\begin{equation}
  \Tr\!\left(
    \left\{
      \nabla\Phi(C_{n+1})+\frac12A_n
      +\frac1h\bigl(\nabla\Phi(C_{n+1})-\nabla\Phi(C_n)\bigr)
    \right\}(C-C_{n+1})
  \right)\geq0.
  \label{eq:app:nonconvex-covariance-vi}
\end{equation}
Set \(C=C_n\) and \(\Delta C_n=C_{n+1}-C_n\).  The Bregman identities
\begin{align*}
  \Tr\!\left(\nabla\Phi(C_{n+1})\Delta C_n\right)
  &=\Phi(C_{n+1})-\Phi(C_n)+D_C^-(C_n,C_{n+1}),\\
  \Tr\!\left(
    \{\nabla\Phi(C_{n+1})-\nabla\Phi(C_n)\}\Delta C_n
  \right)
  &=D_C^-(C_n,C_{n+1})+D_C^+(C_{n+1},C_n)
\end{align*}
therefore give the stronger estimate
\begin{align}
  &\Phi(C_{n+1})-\Phi(C_n)
   +\frac12\Tr\!\left(A_n(C_{n+1}-C_n)\right)
  \notag\\
  &\quad\leq
  -\left(1+\frac1h\right)D_C^-(C_n,C_{n+1})
  -\frac1hD_C^+(C_{n+1},C_n)
  \notag\\
  &\quad\leq
  -\frac1h\left\{
    D_C^-(C_n,C_{n+1})+D_C^+(C_{n+1},C_n)
  \right\}.
  \label{eq:app:nonconvex-covariance-three-point}
\end{align}
Apply \Cref{lem:app:nonconvex-relative-smoothness} with
\(a=a_n\), \(a'=a_{n+1}\), and combine
\eqref{eq:app:nonconvex-mean-optimality}--
\eqref{eq:app:nonconvex-covariance-three-point}.  With
\(L=2\beta\lambda_+\),
\begin{align*}
  \calE(a_{n+1})-\calE(a_n)
  \leq{}&
  -\left(\frac2h-L\right)D_m(a_{n+1},a_n)
  -\frac1hD_C^-(C_n,C_{n+1})\\
  &-\left(\frac1h-L\right)D_C^+(C_{n+1},C_n).
\end{align*}
Because \(hL\leq1\), the final term is nonpositive and
\(2/h-L\geq1/h\).  Hence
\[
  \calE(a_{n+1})
  \leq
  \calE(a_n)-\frac1h\left\{
    \frac12\norm{m_{n+1}-m_n}_{C_n^{-1}}^2
    +D_C^-(C_n,C_{n+1})
  \right\},
\]
which is exactly \eqref{eq:nonconvex-clipped-descent}.  Summation proves the
running-minimum estimate.

Finally, \(\mathsf S_n=0\) if and only if
\((m_{n+1},C_{n+1})=(m_n,C_n)\).  At such a fixed point, mean optimality gives
\(b_n=0\), while \eqref{eq:app:nonconvex-covariance-vi} reduces to
\[
  \Tr\!\left(
    \frac12(A_n-C_n^{-1})(C-C_n)
  \right)\geq0
  \qquad\text{for every feasible }C.
\]
These are exactly the first-order conditions for \(\calE\) on the covariance
box.  Thus \(\mathsf S_n\) is a constrained fixed-point residual, not an
unconstrained Fisher--Rao gradient norm.  Moreover, the full successive
reverse KL is
\[
  \KL\!\left(
    \N(m_n,C_n)\,\middle\|\,\N(m_{n+1},C_{n+1})
  \right)
  =D_C^-(C_n,C_{n+1})
   +\frac12\norm{m_{n+1}-m_n}_{C_{n+1}^{-1}}^2.
\]
Its mean block uses \(C_{n+1}^{-1}\), rather than \(C_n^{-1}\), which is why
the comparison in \Cref{cor:nonconvex-full-kl} incurs the condition-number
factor \(\lambda_+/\lambda_-\).
